% TODO make hw stylesheet
\documentclass[11pt]{article}
\usepackage[letterpaper]{geometry}
\usepackage[usenames,dvipsnames,svgnames]{xcolor}
\usepackage[shortlabels]{enumitem}
\usepackage[framemethod=TikZ]{mdframed}
\usepackage{amsmath,amssymb,amsthm}
\usepackage[colorlinks]{hyperref}
\usepackage{multicol}
\usepackage{tabularx}

\hypersetup{
	colorlinks=true,
	linkcolor=blue,
	filecolor=magenta,
	urlcolor=magenta,
	pdftitle={MAT 324 HW}
}

\usepackage{mathtools}
\usepackage{thmtools}
\usepackage[textsize=footnotesize]{todonotes}
\usepackage{titling}

\mdfdefinestyle{mdbluebox}{roundcorner=10pt,innerbottommargin=9pt,
	linecolor=blue,backgroundcolor=TealBlue!5,}
\declaretheoremstyle[headfont=\sffamily\bfseries\color{MidnightBlue},
mdframed={style=mdbluebox},]{thmbluebox}

\mdfdefinestyle{mdbloobox}{frametitlefont=\bfseries,innerbottommargin=8pt,
	nobreak=true,backgroundcolor=TealBlue!5,linecolor=blue,}
\declaretheoremstyle[headfont=\bfseries\color{MidnightBlue},
mdframed={style=mdbloobox},headpunct={\\[3pt]},postheadspace=0pt,]{thmbloobox}

\mdfdefinestyle{mdredbox}{frametitlefont=\bfseries,innerbottommargin=8pt,
	nobreak=true,backgroundcolor=Salmon!5,linecolor=RawSienna,}
\declaretheoremstyle[headfont=\bfseries\color{RawSienna},
mdframed={style=mdredbox},headpunct={\\[3pt]},postheadspace=0pt,]{thmredbox}

\mdfdefinestyle{mdgreenbox}{linecolor=ForestGreen,backgroundcolor=ForestGreen!5,
	linewidth=2pt,rightline=false,leftline=true,topline=false,bottomline=false,}
\declaretheoremstyle[headfont=\bfseries\sffamily\color{ForestGreen!70!black},
mdframed={style=mdgreenbox},headpunct={ --- },]{thmgreenbox}

\mdfdefinestyle{mdorangebox}{linecolor=RedOrange,backgroundcolor=RedOrange!5,
	linewidth=2pt,rightline=false,leftline=true,topline=false,bottomline=false,}
\declaretheoremstyle[headfont=\bfseries\sffamily\color{RedOrange!70!black},
mdframed={style=mdorangebox},headpunct={ --- },]{thmorangebox}

\mdfdefinestyle{mdblackbox}{linecolor=black,backgroundcolor=RedViolet!5!gray!5,
	linewidth=3pt,rightline=false,leftline=true,topline=false,bottomline=false,}
\declaretheoremstyle[mdframed={style=mdblackbox}]{thmblackbox}

\declaretheoremstyle[spaceabove=6pt,spacebelow=6pt]{basehead}

\declaretheorem[style=basehead,name=Problem]{problem}
\declaretheorem[style=thmredbox,name=Problem,numbered=no]{problem*}
\declaretheorem[style=thmbloobox,name=Setup]{setup} % for config geo
\declaretheorem[style=thmredbox,name=Example,sibling=problem]{example}
\declaretheorem[style=thmredbox,name=$\bigstar$ Problem,sibling=problem]{niceprob}
\declaretheorem[style=thmbluebox,name=Theorem,sibling=problem]{theorem}
\declaretheorem[style=thmbluebox,name=Corollary,sibling=theorem]{corollary}
\declaretheorem[style=thmbluebox,name=Proposition,sibling=theorem]{prop}
\declaretheorem[style=thmbluebox,name=Lemma,numberwithin=problem]{lemma}
\declaretheorem[style=thmbluebox,name=Theorem,numbered=no]{theorem*}
\declaretheorem[style=thmbluebox,name=Lemma,numbered=no]{lemma*}
\declaretheorem[style=thmgreenbox,name=Claim,numberwithin=problem]{claim}
\declaretheorem[style=thmgreenbox,name=Claim,numbered=no]{claim*}
\declaretheorem[style=thmblackbox,name=Remark,numberwithin=problem]{remark}
\declaretheorem[style=thmblackbox,name=Remark,numbered=no]{remark*}
\declaretheorem[style=thmgreenbox,name=Definition,sibling=theorem]{definition}
\declaretheorem[style=thmgreenbox,name=Definition,numbered=no]{definition*}

\addtolength{\textheight}{3.14cm}
\providecommand{\ol}{\overline}
\providecommand{\ceil}[1]{\left\lceil #1 \right\rceil}
\providecommand{\floor}[1]{\left\lfloor #1 \right\rfloor}
\newcommand{\dd}{\mathrm{d}}
\providecommand{\ul}{\underline}
\providecommand{\eps}{\varepsilon}
\providecommand{\half}{\frac{1}{2}}
\providecommand{\dang}{\measuredangle} %% Directed angle
\providecommand{\CC}{\mathbb C}
\providecommand{\FF}{\mathbb F}
\providecommand{\NN}{\mathbb N}
\providecommand{\QQ}{\mathbb Q}
\providecommand{\RR}{\mathbb R}
\providecommand{\ZZ}{\mathbb Z}
\providecommand{\dg}{^\circ}
\providecommand{\ii}{\item}
\providecommand{\setdiff}{\smallsetminus}
\providecommand{\alert}[1]{{\sffamily\textbf{\textcolor{blue}{#1}}}}
\providecommand{\inv}{^{-1}}
\providecommand{\mb}[1]{\mathbf{#1}}

\reversemarginpar

\newenvironment{soln}{\begin{proof}[Solution]}{\end{proof}}

\providecommand{\pow}{\operatorname{Pow}}
\providecommand{\vocab}[1]{{\textbf{\textcolor{ForestGreen}{#1}}}}
\providecommand{\lcm}{\operatorname{lcm}}
\providecommand{\abs}[1]{\left|#1\right|}

\usepackage{asymptote}
\begin{asydef}
	defaultpen(fontsize(10pt));
	unitsize(1cm);
	size(8cm); // set a reasonable default
	usepackage("amsmath");
	usepackage("amssymb");
	settings.tex="pdflatex";
	settings.outformat="pdf";
	// Replacement for olympiad+cse5 which is not standard
	import geometry;
	// recalibrate fill and filldraw for conics
	void filldraw(picture pic = currentpicture, conic g, pen fillpen=defaultpen, pen drawpen=defaultpen)
	{ filldraw(pic, (path) g, fillpen, drawpen); }
	void fill(picture pic = currentpicture, conic g, pen p=defaultpen)
	{ filldraw(pic, (path) g, p); }
	// some geometry
	pair foot(pair P, pair A, pair B) { return foot(triangle(A,B,P).VC); }
	pair orthocenter(pair A, pair B, pair C) { return orthocentercenter(A,B,C); }
	pair centroid(pair A, pair B, pair C) { return (A+B+C)/3; }
	// cse5 abbreviations
	path CP(pair P, pair A) { return circle(P, abs(A-P)); }
	path CR(pair P, real r) { return circle(P, r); }
	pair IP(path p, path q) { return intersectionpoints(p,q)[0]; }
	pair OP(path p, path q) { return intersectionpoints(p,q)[1]; }
	path Line(pair A, pair B, real a=0.6, real b=a) { return (a*(A-B)+A)--(b*(B-A)+B); }
	// cse5 more useful functions
	picture CC() {
		picture p=rotate(0)*currentpicture;
		currentpicture.erase();
		return p;
	}
	pair MP(Label s, pair A, pair B = plain.S, pen p = defaultpen) {
		Label L = s;
		L.s = "$"+s.s+"$";
		label(L, A, B, p);
		return A;
	}
	pair Drawing(Label s = "", pair A, pair B = plain.S, pen p = defaultpen) {
		dot(MP(s, A, B, p), p);
		return A;
	}
	path Drawing(path g, pen p = defaultpen, arrowbar ar = None) {
		draw(g, p, ar);
		return g;
	}
\end{asydef}

\usepackage{mathrsfs}

\setlength{\droptitle}{-8ex}
\pretitle{\begin{flushright}\Large\bfseries}
\posttitle{\par\end{flushright}}
\preauthor{\begin{flushright}}
\postauthor{\end{flushright}}
\predate{\begin{flushright}}
\postdate{\end{flushright}}

\title{\vspace{-2em}MAT 324 HW01}
\author{\large Aditya Pahuja}
\date{\large Due 09/04}
\begin{document}
\maketitle

\begin{problem}
    Let $\mathcal A$ be a collection of closed subsets of $\RR^n$
    such that $\bigcap_{F\in\mathcal A} F=\varnothing$.
    \begin{enumerate}[(a)]
        \ii Assume in addition that $\mathcal A$ contains a bounded subset $F_0$.
	Show that $\mathcal A$ contains a finite subcollection $\mathcal A'$
	so that $\bigcap_{F\in\mathcal A'} = \varnothing$.
	\ii Give an example showing that the above conclusion need not hold
	if $\mathcal A$ does not contain a bounded set.
    \end{enumerate}
\end{problem}

\begin{soln}
    (a) Since the intersection of the sets in $\mathcal A$ is empty,
    we can deduce by de Morgan's law that
    \begin{equation*}
	\bigcup_{F\in\mathcal A}\RR^n\setminus F = \RR^n.
    \end{equation*}
    Since each $F$ in $\mathcal A$ is closed,
    the complements $\RR^n\setminus F$ are open,
    so $\{\RR^n\setminus F : F\in\mathcal A\}$ is an open cover of $F_0$,
    which is compact because it is closed and bounded.
    This means there is a finite subcollection $\{F_1,\cdots,F_n\}\subseteq\mathcal A$
    such that
    \begin{equation*}
	F_0\subseteq \bigcup_{i=1}^n \RR^n\setminus F_i
	\iff
	\RR^n\setminus F_0\supseteq \bigcap_{i=1}^n F_i.
    \end{equation*}
    In particular, the intersection of the $F_i$ is disjoint from $F_0$, so
    $\bigcap_{i=0}^n F_i = \varnothing$;
    thus, $\{F_0,\dots, F_n\}$ is the desired subcollection of $\mathcal A$.

    (b) Let $F_n = [n,\infty)$ be a closed subset of $\RR$
    and let $\mathcal A = \{F_n : n\in\NN\}$; we show that $\mathcal A$
    has the desired property.
    The intersection of all the sets in $\mathcal A$ is empty
    because given $x\in\RR$, we can find $n$ such that $x\notin F_n$,
    namely any $n > x$. However, no finite subcollection of $\mathcal A$
    has empty intersection because, given $a_1,\dots,a_k\in\NN$, we have
    \begin{equation*}
	\bigcap_{i=1}^k F_{a_k} =
	\bigcap_{i=1}^k [a_k, \infty) =
	[\max(a_1, \dots, a_k), \infty)\ne\varnothing.\qedhere
    \end{equation*}
\end{soln}

\newpage

\begin{problem}
    Suppose $a,b\in\RR$ with $a<b$ and $f\colon[a,b]\to\RR$ is a bounded function.
    \begin{enumerate}[(a)]
        \ii Let $\eps > 0$. Show that there exists $\delta > 0$ such that
	\begin{equation*}
	    L(f,\mathcal P) > L(f,[a,b]) - \eps
	    \quad\text{and}\quad
	    U(f,\mathcal P) < U(f,[a,b]) + \eps
	\end{equation*}
	for every partition $\mathcal P$ of $[a,b]$
	with mesh size $\delta(\mathcal P) < \delta$.
	\ii Let $\mathcal P_n$ be a sequence of partitions of $[a,b]$
	such that $\delta(\mathcal P_n)\to 0$. Show that
	\begin{equation*}
	    L(f,[a,b]) = \lim_{n\to\infty} L(f,\mathcal P_n)
	    \quad\text{and}\quad
	    U(f,[a,b]) = \lim_{n\to\infty} U(f,\mathcal P_n).
	\end{equation*}
    \end{enumerate}
\end{problem}

\begin{soln}
    (a) By definition, $L(f,[a,b])$ is the supremum of
    $L(f,\mathcal P)$ over all partitions $\mathcal P$ of $[a,b]$.
    Therefore, for every $\eps > 0$, there exists
    a partition $\mathcal P_{\eps/2} = (x_0, x_1, \dots, x_m)$ such that
    \begin{equation*}
	L(f,\mathcal P_{\eps/2}) > L(f, [a,b]) - \frac\eps2.
    \end{equation*}
    Given an arbitrary partition $\mathcal P = (y_0, y_1, \dots, y_n)$ with
    \begin{equation*}
	\delta(P) \coloneq \min_{i=1,\dots,n} y_i - y_{i-1} <
	\delta \coloneq \frac{\eps}{2M(m+1)},
    \end{equation*}
    we will compare $L(f,\mathcal P)$ with $L(f,\mathcal P')$,
    where $\mathcal P'$ is the common refinement
    $\mathcal P\cup\mathcal P_{\eps/2}$ of $\mathcal P$ and $\mathcal P_{\eps/2}$
    (i.e. merge the two partitions to get $\mathcal P'$).
    Keep in mind that $L(f,\mathcal P')$ is at least as large as
    both $L(f,\mathcal P)$ and $L(f,\mathcal P_{\eps/2})$.
    To do the comparison, observe that at most $m+1$ of the intervals
    $[y_{i-1}, y_i]$ contain some $x_j$, since there are only $m+1$ elements
    in the $\mathcal P_{\eps/2}$ list.
    Clearly, these intervals are the ones that account for
    the difference between $L(f,\mathcal P)$ and $L(f,\mathcal P')$,
    since if $y_{i-1}$ and $y_i$ are consecutive in both $\mathcal P$ and $\mathcal P'$
    then the interval $[y_{i-1},y_i]$ has the same contribution of
    \begin{equation*}
	(y_i - y_{i-1})\inf_{x\in[y_{i-1},y_i]} f(x)
    \end{equation*}
    to both $L(f,\mathcal P)$ and $L(f,\mathcal P')$.
    Within one of these intervals $[y_{i-1}, y_i]$ that contain some $x_j$s,
    we see that the total contribution of this interval to $L(f,\mathcal P')$
    is bounded above by
    \begin{equation*}
	(y_i - y_{i-1})\sup_{x\in[y_{i-1}, y_i]} f(x)
    \end{equation*}
    because the infimum of $f$ over each of the subintervals of $[y_{i-1}, y_i]$
    is trivially bounded above by the suprerum of $f$ over $[y_{i-1}, y_i]$.
    Since $f$ is bounded there exists a positive $M$ such that
    $\abs{f(x)} < M/2$ for all $x$.
    The contribution of the interval $[y_{i-1}, y_i]$ to
    the increase between $L(f,\mathcal P)$ and $L(f,\mathcal P')$
    can then be bounded above as
    \begin{equation*}
	(y_i - y_{i-1})\left(\sup_{x\in[y_{i-1}, y_i]} f(x) -
	\inf_{x\in[y_{i-1}, y_i]}\right) <
	\delta\cdot M.
    \end{equation*}
    There are up to $m+1$ such intervals $[y_{i-1}, y_i]$,
    meaning $L(f, \mathcal P') - L(f, \mathcal P)$ can be bounded above by
    $M(m+1)\delta = \frac{\eps}2$. Therefore,
    \begin{equation*}
        L(f,\mathcal P) > L(f,\mathcal P') - \frac\eps2
	\ge L(f, \mathcal P_{\eps/2}) - \frac\eps2
	> L(f,[a,b]) - \eps,
    \end{equation*}
    which is what we wanted.

    For the upper Riemann integral, we can use the fact that
    $U(f,[a,b]) = -L(-f,[a,b])$ and $U(f,\mathcal P) = -L(-f,\mathcal P)$
    for any partition $\mathcal P$, so, given $\delta$ such that
    $\delta(P) < \delta$ implies $L(-f,\mathcal P) > L(-f,[a,b]) - \eps$,
    \begin{equation*}
	U(f,\mathcal P) = -L(-f,\mathcal P) < -L(-f, [a,b]) - \eps
	= U(f,[a,b]) - \eps.
    \end{equation*}
    Since we want one $\delta$ to satisfy both the condition on $L$ and on $U$,
    we can just take $\delta$ to be the minimum of the two $\delta$s obtained
    for $L$ and $U$ above.

    (b) Given $\eps = \frac 1m$ for some positive integer $m$,
    part (a) posits the existence of a $\delta_m > 0$ such that
    a partition $\mathcal P$ with $\delta(\mathcal P) < \delta_m$ satisfies
    \begin{equation*}
	L(f,\mathcal P) > L(f,[a,b]) - \frac 1m
	\quad\text{and}\quad
	U(f,\mathcal P) < U(f,[a,b]) + \frac 1m.
    \end{equation*}
    Since $\delta(\mathcal P_n)\to 0$, for each $m$, there is an $N_m$
    such that $n > N_m$ implies $\delta(\mathcal P_n) < \delta_m$.
    Thus for each positive integer $m$,
    there exists an $N_m$ such that $n > N_m$ implies
    \begin{align*}
	L(f,[a,b]) \ge L(f,\mathcal P_n) > L(f,[a,b]) - \frac 1m
	&\iff \abs{L(f,[a,b]) - L(f,\mathcal P_n)} < \frac 1m\\
	U(f,[a,b]) \le U(f,\mathcal P_n) < U(f,[a,b]) + \frac 1m
	&\iff \abs{U(f,[a,b]) - U(f,\mathcal P_n)} < \frac 1m.
    \end{align*}
    This is exactly what is meant by the statements
    \begin{equation*}
	L(f,[a,b]) = \lim_{n\to\infty} L(f,\mathcal P_n)
	\quad\text{and}\quad
	U(f,[a,b]) = \lim_{n\to\infty} U(f,\mathcal P_n),
    \end{equation*}
    so we are done.
\end{soln}

\pagebreak

\begin{problem}
    Show that the function $f\colon\RR\to\RR$ given by
    \begin{equation*}
	f(x) = \sum_{n=0}^\infty 2^{-n}\sin(2^nx)
    \end{equation*}
    is well-defined, bounded, and Riemann integrable on $[0,\pi]$.
    Find $\int_0^\pi f(x) \dd x$.
\end{problem}

\begin{soln}
    Since $f_n(x) = 2^{-n}\sin(2^nx)$ is bounded in absolute value by $2^{-n}$
    for all real $x$ and $\sum_{n=0}^\infty 2^{-n} = 2 < \infty$,
    the Weierstass $M$-test tells us that the sum $\sum_{n=0}^\infty f_n(x)$
    converges uniformly on $\RR$.

    We can see that $f$ is bounded because for any nonnegative integer $N$
    the triangle inequality lets us bound the partial sums
    \begin{equation*}
	\abs{\sum_{n=0}^N f_n(x)} \le
	\sum_{n=0}^N\abs{f_n(x)} \le
	\sum_{n=0}^\infty 2^{-n} = 2
    \end{equation*}
    so that in the limit as $N$ grows arbitrarily large we get
    \begin{equation*}
	\abs{f(x)} = \abs{\sum_{n=0}^\infty f_n(x)} \le 2.
    \end{equation*}

    Finally, since $f$ is defined as a uniformly converging
    sum of continuous functions on $\RR$,
    $f$ is itself continuous on $\RR$, hence Riemann integrable on $[0,\pi]$.
    Moreover, the Riemann integral commutes with uniform limits, so
    \begin{align*}
	\int_0^\pi f(x)\dd x
	&= \sum_{n=0}^\infty \int_0^{\pi} 2^{-n}\sin(2^nx)\dd x\\
	&= \sum_{n=0}^\infty 4^{-n}(\cos(2^n\cdot 0) - \cos(2^n\cdot\pi))\\
	&= \cos(0) - \cos(\pi) + \sum_{n=1}^\infty 4^{-n}(\cos(0)-\cos(2^n\pi))\\
	&= 2 + \sum_{n=1}^\infty 4^{-n}(1-1) = \boxed 2\qedhere
    \end{align*}
\end{soln}

\pagebreak

\begin{problem}
    Axler's Definition 2.2 of outer measure on $\RR$ uses \emph{open} intervals.
    Show that the resulting outer measure would not change if \emph{open} intervals
    were replaced by \emph{closed}, or by \emph{closed-open}, or by \emph{open-closed},
    or by \emph{any}, i.e.
    \begin{equation*}
	|A|_{()} = |A|_{[]} = |A|_{[)} = |A|_{(]} = |A|_{\text{any}}
	\qquad \forall A\subseteq\RR.
    \end{equation*}
\end{problem}

\begin{soln}
    Given $A\subseteq\RR$, define
    \begin{equation*}
	\mathcal L_A^{T} = \left\{ \sum_{j=1}^\infty \ell(I_j) :
	I_1, I_2, \dots \text{ are bounded intervals of type $T$ such that }
        A\subseteq\bigcup_{j=1}^\infty I_j\right\},
    \end{equation*}
    where $T$ indicates the status of the intervals' endpoints
    (open, closed, etc.); thus for example
    \begin{equation*}
	\abs{A}_{[)} = \inf\mathcal L_A^{[)}
	\quad\text{and}\quad
	\abs{A}_{\text{any}} = \inf\mathcal L_A^{\text{any}}.
    \end{equation*}
    We stipulate that the intervals $I_j$ are bounded
    because $\infty$ is trivially in each $\mathcal L_A^\bullet$ anyway,
    so we don't need to concern ourselves with intervals of infinite length.
    
    \begin{claim}
        \label{claim:om-ineqs}
	For any $A\subseteq\RR$,
	\begin{equation*}
	    \mathcal L_A^{()}\subseteq
	    \mathcal L_A^{\bullet}\subseteq\mathcal L_A^{\text{any}}.
	\end{equation*}
    \end{claim}

    \begin{proof}
	Obviously $\mathcal L_A^{\text{any}}$ contains
	all the other $\mathcal L_A^\bullet$.
	For the other inclusion, we note that the open interval $(a,b)$
	is contained by each of $(a,b]$, $[a,b)$, and $[a,b]$, so
	if $A$ is covered by a sequence of bounded intervals
	$(a_1, b_1)$, $(a_2, b_2)$, \dots, then
	it is covered by each of the sequences of intervals resulting from
	closing one or both of the endpoints of all the intervals
	(i.e. the sequences $(a_i, b_i]$, $[a_i, b_i)$, and $[a_i, b_i]$);
	therefore if $\sum_{j=1}^\infty (b_j - a_j)$ is in $\mathcal L_A^{()}$
	then it is contained in the other $\mathcal L_A^{\bullet}$ as well.
    \end{proof}
    
    Taking infima gives the inequality
    \begin{equation*}
	\abs{A}_{\text{any}}\le \abs{A}_{\bullet}\le \abs{A}_{()}
    \end{equation*}
    for all $A\subseteq\RR$.
    To finish the solution, we will prove the reverse inequality
    \begin{equation*}
	\abs{A}_{()}\le \abs{A}_{\text{any}}.
    \end{equation*}
    Let $I_1$, $I_2$, \dots be a sequence of (not necessarily open) intervals
    whose union contains $A$. Then, if $I_j$ has endpoints $a_j < b_j$,
    we can see that
    $I_j\subseteq I'_j\coloneq(a_j-\eps\cdot 2^{-j}, b_j+\eps\cdot 2^{-j})$
    for any $\eps>0$, so the $I'_j$ are a covering of $A$ by open intervals.
    Moreover,
    \begin{equation*}
	\underbrace{\sum_{j=1}^\infty (b_j-a_j)}_{\text{total length of $I_j$s}}
	\le \underbrace{\sum_{j=1}^\infty (b_j-a_j + \eps\cdot 2^{-j+1})}_
	{{\text{total length of $I'_j$s}}}
	= \sum_{j=1}^\infty (b_j-a_j) + 2\eps
	= \sum_{j=1}^\infty \ell(I_j) + 2\eps
    \end{equation*}
    so $\abs{A}_{()}\le \sum_{j=1}^\infty\ell(I_j) + 2\eps$
    for all $\eps > 0$, which implies that it is smaller than
    $\sum_{j=1}^\infty\ell(I_j)$. Taking the infimum over
    all choices of the $I_j$s gives $\abs{A}_{()}\le\abs{A}_{\text{any}}$.
\end{soln}

\pagebreak

\begin{problem}
    Show that the set
    \begin{equation*}
	A = \left\{ x\in[-1,1] : |\sin(nx)| \le \frac 1n
	\text{ for all } n\in\NN \right\}
    \end{equation*}
    has (outer) measure zero.
\end{problem}

\begin{soln}
    Let $A_n = \{ x\in[-1,1] : |\sin(nx)| \le \frac 1n\}$, so
    $A = \bigcap_{n=1}^\infty A_n$. This means that $|A| \le |A_n|$
    for all $n$, so it suffices to prove that $\lim_{n\to\infty} |A_n| = 0$.
    To do this, we first rewrite $A_n$ as
    \begin{equation*}
	A_n = \left\{ x\in[-1,1] :
	-\arcsin\frac1n \le nx\bmod\pi \le\arcsin\frac 1n \right\}
    \end{equation*}
    where $nx\bmod\pi$ is the unique real number $r\in[0,\pi]$
    for which there exists an integer $q$ such that $nx = q\pi + r$.
    Equivalently,
    \begin{equation*}
	A_n = \left\{ x\in[-1,1] :
	\frac{-\arcsin\frac1n}n \le x\bmod\frac\pi n \le\frac{\arcsin\frac 1n}n
	\right\}.
    \end{equation*}
    If $x\in[-1,1]$ then $x + \frac\pi n\cdot n > 1$
    and $x - \frac\pi n\cdot n < -1$, so
    each value of $x\bmod\frac\pi n$ corresponds to
    at most $2n$ points in $A_n$. This means that $A_n$ is
    the union of up to $2n$ disjoint closed intervals,
    each with length at most $\frac{2\arcsin\frac 1n}n$, so
    \begin{equation*}
	|A_n|\le \frac{2\arcsin\frac 1n}n \cdot 2n
    \end{equation*}
    Taking the limit as $n$ goes to infinity, we have
    \begin{equation*}
	\abs A \le \lim_{n\to\infty} |A_n| \le
	\lim_{n\to\infty} \frac{2\arcsin\frac 1n}{n} \cdot 2n =
	\lim_{n\to\infty} 4\arcsin\frac 1n =
	4\arcsin(0) = 0
    \end{equation*}
    using the fact that $\arcsin$ is continuous at $0$,
    so $\abs A$ is indeed zero.
\end{soln}

\pagebreak

\begin{problem}
    Let $A\subseteq[0,1]$ be a subset of (outer) measure zero. Show that
    \begin{equation*}
	B = \{ x^2 : x\in A\}
    \end{equation*}
    is also of (outer) measure zero. Is the conclusion still true if $A\subseteq\RR$?
\end{problem}

\begin{soln}
    We prove both parts via the following claim:
    \begin{claim}
        \label{claim:square-meas-zero-interval}
	Let $A$ be a subset of the interval $[a,b]$ for which $0\notin(a,b)$.
	If $A$ has outer measure zero, so does $B = \{x^2 : x\in A\}$.
    \end{claim}
    \begin{proof}
	Let $\eps$ be a positive real number smaller than.
	We know by definition that there exists a sequence of open intervals
	$I_1$, $I_2$, $I_3$, \dots whose union contains $A$ such that
	$\sum_{j=1}^\infty \ell(I_j) < \frac{\eps}{2\max(|a|,|b|)}$.
	Define $I'_j = I_j\cap [a,b]$, so that each interval is contained in $[a,b]$,
	and let $I'_j = (a_j, b_j)$ for real numbers $a_j<b_j$. Then we have
	\begin{equation*}
	    \sum_{j=1}^\infty b_j-a_j =
	    \sum_{j=1}^\infty \ell(I'_j)\le
	    \sum_{j=1}^\infty \ell(I_j) < \frac{\eps}{2\max(|a|,|b|)}.
	\end{equation*}
	Since $x\mapsto x^2$ is monotonic on $[a,b]$ (increasing if $0\le a<b$
	and decreasing if $a<b\le 0$), the set of intervals
	$\{(\min(a_j^2,b_j^2), \max(a_j^2,b_j^2)) : j\in\NN\}$
	covers $B$, and
	\begin{align*}
	    \sum_{j=1}^\infty |b_j^2-a_j^2| &=
	    \sum_{j=1}^\infty (b_j-a_j)\cdot|b_j+a_j|\\
	    &\le \sum_{j=1}^\infty (b_j-a_j)\cdot(|b_j|+|a_j)\\
	    &\le \sum_{j=1}^\infty (b_j-a_j)\cdot2\max(|a|,|b|)\\
	    &< 2\max(|a|,|b|)\cdot \frac{\eps}{2\max(|a|,|b|)} = \eps.
	\end{align*}
	Thus the outer measure of $B$ is bounded above by $\eps$ for all $\eps$,
	i.e. is zero. (Here we are using the result of Problem 4
	so that some of the intervals in the covering $I'_j$ of $A$
	may be half-open instead of open.)
    \end{proof}

    This immediately solves the first part of the problem with $a=0$, $b=1$.
    For the second part of the problem, we can write
    $\RR = \bigcup_{n\in\ZZ}[n,n+1]$, noting that $0$ does not lie in
    the interior of any of these intervals.
    If $A\subseteq\RR$ has measure zero, then by monotonicity
    so do its subsets $A\cap[n,n+1]$ for any $n\in\ZZ$;
    and then by the claim, $B_n = \{x^2 : x\in A\cap[n,n+1]\}$
    also has outer measure zero.
    Obviously $B = \bigcup_{n\in\ZZ} B_n$,
    and the union of countably many sets of outer measure zero
    has outer measure zero (by countable subadditivity,
    since the outer measure of the union is bounded by
    the sum of countably many zeroes),
    which implies that $|B| = 0$.
\end{soln}

\pagebreak

\begin{problem}
    Let $C$ be the Cantor set and $\Lambda\colon[0,1]\to[0,1]$ the Cantor function.
    \begin{enumerate}[(a)]
	\ii Suppose $x$ is a rational number in $[0,1]$.
	Explain why $\Lambda(x)$ is rational.
	\ii Suppose $x\in C$ is such that $\Lambda(x)$ is rational.
	Explain why $x$ is rational.
    \end{enumerate}
\end{problem}

\begin{soln}
    We will take for granted the fact that
    a real number $x\in[0,1]$ is rational if and only if,
    for every integer $b\ge 2$, when $x$ is expressed in base $b$ as
    \begin{equation*}
	x = \sum_{n=1}^\infty \frac{d_n}{b^n}
    \end{equation*}
    with digits $0\le d_n<n$, the sequence $(d_n)_{n\ge r}$ is
    eventually periodic, which we say to mean that there exists an $N$
    such that $(d_n)_{n\ge N}$ is a periodic sequence.

    (a) Let $x\in[0,1]$ be rational. If $x\in C$, then
    \begin{equation*}
	x \overset{\text{base 3}}= \sum_{n=1}^\infty \frac{d_n}{3^n}
	\implies
	\Lambda(x) \overset{\text{base 2}}= \sum_{n=1}^\infty \frac{d_n/2}{2^n}.
    \end{equation*}
    where the digits $d_n$ of the base-3 expansion of $x$ are in $\{0,2\}$.
    Since $(d_n)_{n\ge1}$ is eventually periodic (because $x$ is rational),
    so is the sequence $(d_n/2)_{n\ge1}$, so $\Lambda(x)$ is rational.
    On the other hand, if $x\notin C$, then
    let $t$ be the smallest positive integer for which
    the base-3 expansion of $x$ has a digit $d_t = 1$. Then,
    \begin{equation*}
	\Lambda(x) = \sum_{n=1}^{t-1} \frac{d_n/2}{2^n}
    \end{equation*}
    which is rational because it is the sum of finitely many rational numbers.

    (b) As stated in the solution to (a), if $x\in C$ then
    \begin{equation*}
	x \overset{\text{base 3}}= \sum_{n=1}^\infty \frac{d_n}{3^n}
	\implies
	\Lambda(x) \overset{\text{base 2}}= \sum_{n=1}^\infty \frac{d_n/2}{2^n}.
    \end{equation*}
    Therefore if $\Lambda(x)$ is rational for some $x\in C$,
    the sequence $(d_n/2)_{n\ge1}$ is eventually periodic,
    hence so is the sequence $(d_n)_{n\ge1}$,
    implying that $x$ is rational too.
\end{soln}

\end{document}
